%----------------------------------------------------------------------------------------
%   CLASS DEFINITIONS (EMBEDDED RESUME.CLS)
%----------------------------------------------------------------------------------------
\documentclass[9pt,letterpaper]{article} % Font size and paper type

\usepackage[parfill]{parskip} % Remove paragraph indentation
\usepackage{array} % Required for boldface (\bf and \bfseries) tabular columns
\usepackage{ifthen} % Required for ifthenelse statements
\usepackage{charter}
\usepackage{enumitem}
\usepackage{xcolor} 
\usepackage[colorlinks=true, urlcolor=blue]{hyperref}
\usepackage[left=0.5in,top=0.25in,right=0.4in,bottom=0.2in, columnsep=0.5cm]{geometry} 
\input{glyphtounicode}
\usepackage[none]{hyphenat}
\usepackage{multicol}

\pdfgentounicode=1
\pagestyle{empty} % Suppress page numbers

\makeatletter % Begin internal command redefinitions

%----------------------------------------------------------------------------------------
%	HEADINGS COMMANDS: Commands for printing name and address
%----------------------------------------------------------------------------------------

\def \name#1{\def\@name{#1}} % Defines the \name command to set name
\def \@name {} % Sets \@name to empty by default

\def \addressSep {$\diamond$} % Set default address separator to a diamond

% One, two or three address lines can be specified 
\let \@addressone \relax
\let \@addresstwo \relax
\let \@addressthree \relax

% \address command can be used to set the first, second, and third address (last 2 optional)
\def \address #1{
  \@ifundefined{@addresstwo}{
    \def \@addresstwo {#1}
  }{
  \@ifundefined{@addressthree}{
  \def \@addressthree {#1}
  }{
     \def \@addressone {#1}
  }}
}

% \printaddress is used to style an address line (given as input)
\def \printaddress #1{
  \begingroup
    \def \\ {\addressSep\ }
    \centerline{#1}
  \endgroup
  \par
}

% \printname is used to print the name as a page header
\def \printname {
  \begingroup
    \hfil{{\namesize\bf \@name}}\hfil
    \nameskip\break
  \endgroup
}

%----------------------------------------------------------------------------------------
%	PRINT THE HEADING LINES
%----------------------------------------------------------------------------------------

\let\ori@document=\document
\renewcommand{\document}{
  \ori@document  % Begin document
  \printname % Print the name specified with \name
  \@ifundefined{@addressone}{}{ % Print the first address if specified
    \printaddress{\@addressone}}
  \@ifundefined{@addresstwo}{}{ % Print the second address if specified
    \printaddress{\@addresstwo}}
     \@ifundefined{@addressthree}{}{ % Print the third address if specified
    \printaddress{\@addressthree}}
}

%----------------------------------------------------------------------------------------
%	SECTION FORMATTING
%----------------------------------------------------------------------------------------

% Defines the rSection environment for the large sections within the CV
\newenvironment{rSection}[1]{ % 1 input argument - section name
  \sectionskip
 \textbf{\Large #1} % Section title
  \sectionlineskip
  \hrule % Horizontal line
  \begin{list}{}{ % List for each individual item in the section
    \setlength{\leftmargin}{-0em} % Margin within the section
  }
  \item[]
}{
  \end{list}
}

%----------------------------------------------------------------------------------------
%	WORK EXPERIENCE FORMATTING
%----------------------------------------------------------------------------------------

\newenvironment{rSubsection}[4]{ % 4 input arguments - company name, year(s) employed, job title and location
 {\bf #1} \hfill {#2} % Bold company name and date on the right
 \ifthenelse{\equal{#3}{}}{}{ % If the third argument is not specified, don't print the job title and location line
  \\
  {\em #3} \hfill {\em #4} % Italic job title and location
  }\smallskip
  \begin{list}{$\cdot$}{\leftmargin=0em} % \cdot used for bullets, no indentation
   \itemsep -0.5em \vspace{-0.5em} % Compress items in list together for aesthetics
  }{
  \end{list}
  \vspace{0em} % Some space after the list of bullet points
}

% The below commands define the whitespace after certain things in the document - they can be \smallskip, \medskip or \bigskip
\def\namesize{\huge} % Size of the name at the top of the document
\def\addressskip{} % The space between the two address (or phone/email) lines
\def\sectionlineskip{\smallskip} % The space above the horizontal line for each section 
\def\nameskip{\vspace{-0.5em}} % The space after your name at the top
\def\sectionskip{} % The space after the heading section

\makeatother % End internal command redefinitions

%----------------------------------------------------------------------------------------
%   DOCUMENT START
%----------------------------------------------------------------------------------------
\name{Abhishek Nagaraja, (CSPO) \vspace{.2em} \hrule}

\address{Arlington, Texas, USA \textbullet\ work.abhishekn@gmail.com \textbullet\ {\href{https://www.abhishekn.in}{www.abhishekn.in}} \textbullet\ {\href{https://linkedin.com/in/abhisheknagaraja}{linkedin.com/in/abhisheknagaraja}}}

\begin{document}

%----------------------------------------------------------------------------------------
%   SUMMARY
%----------------------------------------------------------------------------------------

 Certified Scrum Product Owner (CSPO) with experience in \textbf{Agile Product Delivery} and \textbf{Backlog Management}. Proven ability to translate complex business requirements into clear \textbf{User Stories} and technical specifications. Expert in facilitating Scrum ceremonies, aligning stakeholders, and leveraging AI/Cloud technologies to accelerate product velocity.


\begin{rSection}{Education}
\begin{rSubsection}{University of Texas at Arlington}{Texas, USA}{\textbf{Master of Science in Computer Science}}{August 2023 - May 2024}
\item \textbf{Master of Science in Engineering Management} \hfill May 2024 - Current
\end{rSubsection}

\begin{rSubsection}{Jawaharlal Nehru Technological University}{Hyderabad, India}{\textbf{Bachelor of Technology in Computer Science and Engineering}}{August 2019 - July 2023}
\item[] \vspace{-1em}
\end{rSubsection}
\end{rSection}

%----------------------------------------------------------------------------------------
%   WORK EXPERIENCE SECTION
%----------------------------------------------------------------------------------------
\begin{rSection}{Work Experience}

\begin{rSubsection}
{Program Manager Intern, University of Texas at Arlington - Texas, USA}{October 2023 - October 2025}{}{}
\item[] \vspace{-0.5em}
\begin{itemize}[left=-1.5em, labelsep=.1cm, labelwidth=.2cm, itemsep=0.0em]
\item Accomplished \textbf{successful product launch} as measured by \textbf{scaling to 5,000+ users and \$120,000 in revenue}, by \textbf{acting as Product Owner} for 'MavMarket' to manage the full product lifecycle.
\item Accomplished \textbf{40\% reduction in overhead} as measured by \textbf{automated departmental workflows}, by \textbf{architecting an AI-powered Standard Operating Procedure (SOP)} to groom the operational backlog and remove bottlenecks.
\item Accomplished \textbf{35\% increase in Net Promoter Score (NPS)} as measured by \textbf{stakeholder feedback loops}, by \textbf{using data analytics} to prioritize improvements and drive user engagement metrics up by \textbf{25\%}.
\item Accomplished \textbf{on-time delivery of objectives} as measured by \textbf{seamless event execution}, by \textbf{facilitating cross-functional collaboration} between marketing, logistics, and vendors.
\end{itemize}
\end{rSubsection}

\begin{rSubsection}
{Business Analyst Intern, NFC Solutions Private Limited - Hyderabad, India}{April 2022 - January 2023}{}{}
\item[] \vspace{-0.5em}
\begin{itemize}[left=-1.5em, labelsep=.1cm, labelwidth=.2cm, itemsep=0.0em]
\item Accomplished \textbf{aligned technical execution} as measured by \textbf{delivery of 4 key projects}, by \textbf{serving as Proxy Product Owner} to translate business needs into \textbf{Product Requirement Documents (PRDs)} and \textbf{API specifications}.
\item Accomplished \textbf{20\% improvement in velocity} as measured by \textbf{development efficiency metrics}, by \textbf{facilitating Sprint Planning, Daily Scrums, and Retrospectives} to remove blockages.
\item Accomplished \textbf{smooth API integrations} as measured by \textbf{successful feature releases}, by \textbf{managing the Product Backlog} for an e-commerce platform and prioritizing features based on customer value.
\end{itemize}
\end{rSubsection}

\begin{rSubsection}
{Founder and Community Lead, e-DAM Community - Hyderabad, India}{March 2021 - January 2025}{}{}
\item[] \vspace{-0.5em}
\begin{itemize}[left=-1.5em, labelsep=.1cm, labelwidth=.2cm, itemsep=0.0em]
\item Accomplished \textbf{community scale} as measured by \textbf{growth to 5,000+ members}, by \textbf{defining the Product Strategy} for a technical platform and releasing iterative features based on user feedback.
\item Accomplished \textbf{high-value content delivery} as measured by \textbf{200+ technical workshops}, by \textbf{managing stakeholder relationships} with corporate partners to align offerings with market demand.
\end{itemize}
\end{rSubsection}

\end{rSection}

%----------------------------------------------------------------------------------------
%   PROJECTS
%----------------------------------------------------------------------------------------
\begin{rSection}{Projects}

\begin{rSubsection}{Thara: Multi-Agent AI Ecosystem}{}{}{}
\item[] \vspace{-0.5em}
\begin{itemize}[left=-1.5em, labelsep=.1cm, labelwidth=.2cm, itemsep=0.0em]
\item Accomplished \textbf{streamlined development workflows} as measured by \textbf{actionable User Stories for 11 components}, by \textbf{defining and prioritizing the Product Backlog} for a complex, scalable AI infrastructure.
\item Accomplished \textbf{efficient cross-functional coordination} as measured by \textbf{on-time sprint delivery}, by \textbf{operating as Product Owner} to facilitate ceremonies across \textbf{Cloud (AWS)} and \textbf{Backend (FastAPI)} engineering workstreams.
\item Accomplished \textbf{99.9\% system uptime} as measured by \textbf{automated validation scripts}, by \textbf{defining clear Acceptance Criteria} and overseeing quality execution.
\end{itemize}
\end{rSubsection}

\begin{rSubsection}{Google Maps: Product Teardown}{}{}{}
\item[] \vspace{-0.5em}
\begin{itemize}[left=-1.5em, labelsep=.1cm, labelwidth=.2cm, itemsep=0.0em]
\item Accomplished \textbf{validated feature prioritization} as measured by \textbf{42 user interviews}, by \textbf{translating qualitative feedback} into prioritized \textbf{Epics} focused on "Menu Accuracy" and "Dietary Filtering."
\item Accomplished \textbf{stakeholder buy-in} as measured by \textbf{approval of high-impact features}, by \textbf{developing a Product Roadmap} for a centralized menu system addressing critical user needs.
\end{itemize}
\end{rSubsection}

\begin{rSubsection}{Clash of Clans: Product Teardown}{}{}{}
\item[] \vspace{-0.5em}
\begin{itemize}[left=-1.5em, labelsep=.1cm, labelwidth=.2cm, itemsep=0.0em]
\item Accomplished \textbf{targeted user growth} as measured by \textbf{5-10\% estimated organic acquisition}, by \textbf{prioritizing feature development} based on quantitative impact analysis of social sharing.
\item Accomplished \textbf{projected retention lift} as measured by \textbf{15-20\% increase in join rates}, by \textbf{defining release requirements} for a "Clan Discovery Engine" aligned with retention goals.
\end{itemize}
\end{rSubsection}

\end{rSection}

%----------------------------------------------------------------------------------------
%   SKILLS
%----------------------------------------------------------------------------------------
\begin{rSection}{Skills and Tools}
\item[] \vspace{-0.5em}
\begin{itemize}[left=-1.5em, labelsep=.1cm, labelwidth=.2cm, itemsep=0.0em]
    \item \textbf{Product Ownership:} Backlog Grooming, User Story Writing, Sprint Planning, Acceptance Criteria, Burndown Charts, Stakeholder Management, Agile/Scrum.
    \item \textbf{Tools:} Jira, Confluence, Azure DevOps, Asana, Notion, Figma, Power BI, Google Analytics.
    \item \textbf{Technical Literacy:} API Integration, REST, SQL, Cloud Architecture (AWS), Python, AI Models (Gemini, GPT-4), CI/CD Concepts.
    \item \textbf{Methodologies:} Scrum, Kanban, Lean-Agile, RICE Prioritization, Design Thinking.
\end{itemize}
\end{rSection}

\begin{rSection}{Certification}
\item[] \vspace{-1em}
\begin{itemize}[left=-1.5em, labelsep=.1cm, labelwidth=.2cm, itemsep=0.0em]
\item \href{https://drive.google.com/file/d/1W_00yFKzw16QAXMxebXWl7uEZn83txCp/view?usp=sharing}{Certified Scrum Product Owner (CSPO)}
\item \href{https://www.coursera.org/account/accomplishments/verify/AZCR52ERTR6U}{Google: Foundation of Project Management}
\item Google: Data Analytics Specialization
\end{itemize}
\end{rSection}

\begin{rSection}{Other Projects}
\begin{rSubsection}{Product-Pal: AI Strategy Specialist}{}{}{}
\item[] \vspace{-0.5em}
\begin{itemize}[left=-1.5em, labelsep=.1cm, labelwidth=.2cm, itemsep=0.0em]
\item Accomplished \textbf{70\% reduction in analysis time} as measured by \textbf{operational efficiency}, by \textbf{defining the Minimum Viable Product (MVP)} for an autonomous AI strategy agent.
\item Accomplished \textbf{reduced API latency} as measured by \textbf{user testing metrics}, by \textbf{validating product assumptions} through rapid prototyping and UI iteration.
\end{itemize}
\end{rSubsection}

\begin{rSubsection}{Split-Wiser: AI-Powered Financial Automation}{}{}{}
\item[] \vspace{-0.5em}
\begin{itemize}[left=-1.5em, labelsep=.1cm, labelwidth=.2cm, itemsep=0.0em]
\item Accomplished \textbf{prototype delivery} as measured by \textbf{completed concept-to-launch cycle}, by \textbf{owning the Product Roadmap} for a financial automation tool using OpenAI APIs.
\end{itemize}
\end{rSubsection}

\end{rSection}

\end{document}
